\section{Results}

\subsection{Internal and Calendar-Aware Evaluation}

The selected participant-disjoint multimodal model combined cough and speech
through a validation-fitted logistic stack. It achieved AUROC 0.897, AUPRC
0.863, balanced accuracy 0.825, and F1 0.752 on 314 test participants. Under the
time-stratified participant protocol, the selected uniform cough--breath--speech
fusion achieved AUROC 0.849, AUPRC 0.783, balanced accuracy 0.783, and F1 0.705
on 431 participants. These two rows describe protocol-specific selected models;
their numerical difference is not a paired estimate of the effect of calendar
stratification.

Across available seed replications, the validation-stacked internal
cough--speech fusion had mean AUROC $0.895\pm0.003$ (four seeds). The selected
early-to-late breath ensemble had mean AUROC $0.691\pm0.007$ (three seeds).
The early-to-late selected row reached AUROC 0.698 but had ECE 0.711, indicating
that the probabilities were poorly aligned with observed outcomes. Again,
these rows involve different final modality combinations and are evidence of
protocol sensitivity across the selected workflow, not a controlled paired
model comparison.

The exploratory reverse temporal run produced AUROC 0.920 but AUPRC 0.029,
F1 0.011, and ECE 0.471. The high AUROC therefore did not imply a useful
operating point under the reversed prevalence and cohort composition.

\subsection{Modality-Matched External Transfer}

Table~\ref{tab:external} reports like-for-like cough transfer. Every
conventional model lost substantial discrimination when moved from Coswara to
COUGHVID. In the independent two-sample bootstrap, Coswara cough predictions
were first averaged within each participant and participants were resampled;
COUGHVID recordings were resampled independently. All four 95\% intervals for
the AUROC difference excluded zero.

\begin{table}[!t]
\centering
\caption{Cough-only internal-to-external transfer. Conventional-model
differences include participant-cluster-correct 95\% bootstrap intervals.}
\label{tab:external}
\begin{threeparttable}
\begin{tabular}{lccc}
\toprule
Model family & Internal AUROC & External AUROC & Difference (95\% CI) \\
\midrule
LightGBM & 0.853 & 0.543 & 0.310 (0.247, 0.366) \\
SVC--RBF & 0.868 & 0.523 & 0.345 (0.291, 0.398) \\
CatBoost & 0.849 & 0.531 & 0.318 (0.256, 0.376) \\
XGBoost & 0.850 & 0.532 & 0.318 (0.262, 0.371) \\
WavLM Base Plus & 0.812 & 0.484 & 0.328 (not estimated) \\
CNN--BiGRU & 0.737 & 0.548 & 0.189 (not estimated) \\
\bottomrule
\end{tabular}
\begin{tablenotes}\footnotesize
\item The external prevalence was 3.42\% (285/8,331). COUGHVID evaluates only
the cough branch. Conventional internal values are participant-level means over
available Coswara cough recordings ($n=316$); external values use 8,331 target
recordings. Deep-model differences are descriptive because comparable
cluster-correct intervals were not estimated for those runs.
\end{tablenotes}
\end{threeparttable}
\end{table}

At required sensitivity of at least 0.90, the conventional external models
retained only approximately 0.11--0.16 specificity and 0.035--0.037 precision.
That precision was close to the 0.034 target prevalence. This means that, at
the chosen screening constraint, a positive prediction provided little gain
over the target base rate.

Target-only Platt and isotonic recalibration changed probability scale and
threshold-dependent metrics but did not materially change AUROC. The result
separates calibration error from ranking failure: threshold adjustment can
alter decisions, but cannot recover discrimination that is absent in the score
ordering.

Figure~\ref{fig:validation-transfer} keeps the protocol-specific selected
systems separate from the modality-matched transfer comparison. This prevents
the 0.897 multimodal internal value from being presented as the source endpoint
of a cough-only external experiment.

\begin{figure}[!t]
\centering
\includegraphics[width=0.99\linewidth]{fig02_validation_and_transfer.pdf}
\caption{Discrimination under the evaluated protocols. (A) Protocol-specific
selected systems; modality combinations differ and the points are descriptive.
(B) Cough-only internal-to-COUGHVID transfer across conventional and learned
representations. Brackets give 95\% bootstrap intervals for conventional-model
AUROC differences.}
\label{fig:validation-transfer}
\end{figure}

\subsection{Metadata Association and Incremental Audio Value}

The full safe metadata model reached AUROC 0.964; symptoms alone reached 0.932,
and demographic plus recording-protocol variables reached 0.914. After full
model retraining with permuted labels, mean AUROC values were 0.503, 0.503, and
0.503, respectively. The negative control argues against a simple leakage or
metric-calculation bug in the metadata pipeline.

Grouped coefficient magnitude and permutation analyses identified recording
year and recording-protocol variables as influential. This demonstrates that
label-predictive collection context exists. It does not by itself prove that
the audio model encoded the same variables, because the metadata and audio
predictors are distinct measurements.

Incremental value was evaluated on candidates selected by validation
performance before test comparison. For the highest-ranked aligned
cough--speech candidate, the full metadata model achieved test AUROC 0.924 and
metadata plus audio achieved 0.935. The difference was 0.012 (95\% bootstrap
interval $-0.033$ to 0.053; paired DeLong $p=0.572$; $n=61$). Against the
symptoms-only model, adding the same audio candidate changed AUROC by 0.063
(95\% bootstrap interval $-0.005$ to 0.149; paired DeLong $p=0.104$). These
small aligned samples do not establish absence of incremental audio value, but
they do not provide precise evidence of an improvement.

\subsection{Mechanism and Robustness Analyses}

Only 110 of the top 800 acoustic features ranked independently in the early
period were also ranked in the late period. The union contained 1,490 features,
giving a Jaccard overlap of 0.074. This sensitivity analysis indicates that the
feature-ranking output was temporally unstable. Because prevalence and case mix
also changed, it neither isolates the cause nor proves that underlying biology
changed.

A source-versus-target classifier based on 500 common features achieved AUROC
0.750. Using the 5th--95th percentile band of source-domain probabilities,
25.2\% of external recordings fell outside the source band. This diagnostic
supports meaningful source-target mismatch and warns that reweighting or
covariance alignment may require extrapolation.

Audio label-shuffle retraining produced AUROC values of 0.537--0.551 across the
three final protocols. These are close enough to chance to support the absence
of gross label leakage, while their finite-sample deviation from exactly 0.5
is retained rather than rounded away. Quality flags, duration, subgroup, and
matched-cohort analyses were treated as sensitivity analyses and are
reported in the supplementary evidence tables.

Figure~\ref{fig:mechanism} collects the principal mechanism and negative-control
results. The panels quantify label-predictive context, imprecise incremental
audio value, temporal feature-ranking instability, and source-target mismatch;
none alone identifies a causal shortcut encoded by the audio model.

\begin{figure}[!t]
\centering
\includegraphics[width=0.99\linewidth]{fig03_mechanism_and_incremental_value.pdf}
\caption{Mechanism and sensitivity evidence. (A) Metadata AUROC before and
after full shuffled-label retraining. (B) AUROC change after adding the
validation-ranked audio candidate to metadata baselines. (C) overlap of early
and late top-800 feature sets. (D) fraction of COUGHVID recordings inside the
source-domain classifier's 5th--95th percentile probability band.}
\label{fig:mechanism}
\end{figure}
